\section{Cauchy--Binet and related formulas}

We work over a commutative ring $K$. We formalize the Cauchy--Binet formula
for determinants of products of non-square matrices, the expansion of
$\det(A+B)$ as a sum over pairs of subsets, and several applications:
formulas for $\det(A+D)$ when $D$ is diagonal, the characteristic
polynomial in terms of principal minors, and the determinant of the Pascal
matrix via LU decomposition.

Throughout, $[n] = \{1,2,\ldots,n\}$ denotes the set of the first $n$
positive integers (in the Lean formalization, indexing is $0$-based).
For any finite set $S$ of integers we write
$\operatorname{sum} S := \sum_{s \in S} s$.

\subsection{Submatrices and minors}

\begin{definition}
\label{def.det.sub}
\lean{AlgebraicCombinatorics.CauchyBinet.submatrixOfFinset}
\leantarget
\leanok
Let $n,m\in\mathbb{N}$. Let $A$ be an $n\times m$-matrix.

Let $U$ be a subset of $[n]$. Let $V$ be a subset of $[m]$.

Writing the two sets $U$ and $V$ as
\[
U=\{u_1,u_2,\ldots,u_p\}
\quad\text{and}\quad
V=\{v_1,v_2,\ldots,v_q\}
\]
with
\[
u_1<u_2<\cdots<u_p
\quad\text{and}\quad
v_1<v_2<\cdots<v_q,
\]
we set
\[
\operatorname{sub}_U^V A:=\left(A_{u_i,v_j}\right)_{1\le i\le p,\ 1\le j\le q}.
\]

Roughly speaking, $\operatorname{sub}_U^V A$ is the matrix obtained from
$A$ by focusing only on the $i$-th rows for $i\in U$ (that is,
removing all the other rows) and only on the $j$-th columns for
$j\in V$ (that is, removing all the other columns).

This matrix $\operatorname{sub}_U^V A$ is called the \emph{submatrix of}
$A$ \emph{obtained by restricting to the $U$-rows
and the $V$-columns}. If this matrix is square (i.e., if $|U|=|V|$), then its determinant
$\det(\operatorname{sub}_U^V A)$ is called a \emph{minor} of
$A$.
\end{definition}

\begin{lemma}
\label{lem.minor.empty}
\lean{AlgebraicCombinatorics.CauchyBinet.minor_empty}
\leanhelper
\leanok
The minor of the empty sets is $1$:
$\det(\operatorname{sub}_{\varnothing}^{\varnothing} A) = 1$.
\end{lemma}

\begin{proof}
\leanok
The $0 \times 0$ matrix has determinant $1$.
This is \texttt{Mathlib.Matrix.det\_isEmpty} from Mathlib.
\end{proof}

\subsection{The Cauchy--Binet formula}

\subsubsection{Helpers for the Cauchy--Binet formula}

\begin{lemma}[Non-injective cancellation]
\label{lem.det.CB.nonsquare}
\lean{AlgebraicCombinatorics.CauchyBinet.det_mul_aux_nonsquare}
\leanhelper
\leanok
Let $A\in K^{n\times m}$, $B\in K^{m\times n}$, and let
$f\colon [n]\to [m]$ be a non-injective function.  Then
\[
\sum_{\sigma\in S_n} (-1)^{\sigma}\,
\prod_{i=1}^n A_{\sigma(i),f(i)}\, B_{f(i),i}
= 0.
\]
\end{lemma}

\begin{proof}
Since $f$ is not injective, there exist $i\ne j$ with $f(i)=f(j)$.
Pairing each $\sigma$ with $\sigma\cdot\operatorname{swap}(i,j)$
gives a sign-reversing involution, so the sum cancels to~$0$.
\end{proof}

\begin{lemma}[Subset factorization]
\label{lem.det.CB.subset_factor}
\lean{AlgebraicCombinatorics.CauchyBinet.sum_over_subset_eq_det_mul}
\leanhelper
\leanok
For a fixed $n$-element subset $S\subseteq [m]$,
\[
\sum_{\tau\in S_n}\sum_{\sigma\in S_n}
(-1)^{\sigma}\prod_{i=1}^n
A_{\sigma(i),S_{\tau(i)}}\,B_{S_{\tau(i)},i}
= \det(\operatorname{cols}_S A)\cdot\det(\operatorname{rows}_S B),
\]
where $S_k$ denotes the $k$-th element of $S$ in sorted order.
\end{lemma}

\begin{proof}
Split the product $\prod_i A_{\sigma(i),e(\tau(i))}\,B_{e(\tau(i)),i}$
into a product of $A$-factors and a product of $B$-factors.
The $B$-factors depend only on~$\tau$ and can be pulled out of the
$\sigma$-sum. The remaining $\sigma$-sum is a Leibniz expansion
for $\det(\operatorname{cols}_S A)$, and summing the $\tau$-terms
gives $\det(\operatorname{rows}_S B)$.
\end{proof}

\begin{theorem}[Cauchy--Binet formula]
\label{thm.det.CB}
\lean{AlgebraicCombinatorics.CauchyBinet.cauchyBinet}
\leantarget
\leanok
Let $n,m\in\mathbb{N}$. Let $A\in K^{n\times m}$ be an $n\times m$-matrix,
and let $B\in K^{m\times n}$ be an $m\times n$-matrix. Then,
\[
\det(AB) = \sum_{\substack{(g_1,g_2,\ldots,g_n)\in [m]^n;\\
g_1<g_2<\cdots<g_n}}
\det(\operatorname{cols}_{g_1,g_2,\ldots,g_n} A)
\cdot \det(\operatorname{rows}_{g_1,g_2,\ldots,g_n} B).
\]
Here, $\operatorname{cols}_{g_1,\ldots,g_n} A$ is the $n\times n$-matrix
obtained from $A$ by keeping only columns $g_1,g_2,\ldots,g_n$, and
$\operatorname{rows}_{g_1,\ldots,g_n} B$ is the $n\times n$-matrix
obtained from $B$ by keeping only rows $g_1,g_2,\ldots,g_n$.

Equivalently, the sum runs over all
$n$-element subsets $S$ of $[m]$:
\[
\det(AB) = \sum_{\substack{S \subseteq [m] \\ |S|=n}}
\det(\operatorname{cols}_S A) \cdot \det(\operatorname{rows}_S B).
\]
\end{theorem}

\begin{proof}

The proof expands $\det(AB)$ via the Leibniz formula:
\[
\det(AB) = \sum_{\sigma\in S_n} (-1)^{\sigma}
\prod_{i=1}^n (AB)_{\sigma(i),i}
= \sum_{\sigma\in S_n} (-1)^{\sigma}
\prod_{i=1}^n \sum_{k=1}^m A_{\sigma(i),k} B_{k,i}.
\]

Expanding the product over $i$ as a sum over functions
$f\colon [n]\to[m]$ and swapping sums gives
\[
\det(AB) = \sum_{f\colon [n]\to [m]}
\sum_{\sigma\in S_n} (-1)^{\sigma}
\prod_{i=1}^n A_{\sigma(i),f(i)} B_{f(i),i}.
\]

When $f$ is not injective, the inner sum vanishes by
Lemma~\ref{lem.det.CB.nonsquare}.

Therefore only injective $f$ contribute. Each injective $f$ has image
$S = \operatorname{im}(f)$ with $|S|=n$. The injective functions with
a given image $S$ are in bijection with permutations of $[n]$.
Partitioning by $S$ and reindexing via this bijection,
the sum transforms into a sum over $n$-element subsets $S$, and
for each $S$ the inner sum factors as
$\det(\operatorname{cols}_S A)\cdot \det(\operatorname{rows}_S B)$
by Lemma~\ref{lem.det.CB.subset_factor}.
\end{proof}

\begin{theorem}[Cauchy--Binet for square matrices]
\label{thm.det.CB.square}
\lean{AlgebraicCombinatorics.CauchyBinet.cauchyBinet_square}
\leanhelper
\leanok
When $m=n$, the Cauchy--Binet formula specializes to
$\det(AB) = \det A \cdot \det B$.
\end{theorem}

\begin{proof}
\leanok
This is Mathlib's \texttt{Matrix.det\_mul}.
\end{proof}

\subsection{The formula for $\det(A+B)$}

\subsubsection{Basic definitions}

\begin{definition}[Subset sum and permutation sets]
\label{def.finsetSumFin}
\lean{AlgebraicCombinatorics.CauchyBinet.finsetSumFin, AlgebraicCombinatorics.CauchyBinet.permsMapping, AlgebraicCombinatorics.CauchyBinet.imageFinset}
\leanhelper
\leanok
For a subset $P\subseteq [n]$, define
$\operatorname{sum}(P) = \sum_{i\in P} i$.

For subsets $P,Q\subseteq [n]$, define
$\operatorname{Perm}(P,Q) = \{\sigma\in S_n : \sigma(P)=Q\}$
where $\sigma(P) = \{\sigma(i) \mid i \in P\}$
denotes the image of $P$ under $\sigma$.
\end{definition}

\begin{lemma}
\label{lem.permsMapping.empty}
\lean{AlgebraicCombinatorics.CauchyBinet.permsMapping_empty_of_card_ne}
\leanhelper
\leanok
If $|P| \ne |Q|$, then $\operatorname{Perm}(P,Q) = \varnothing$.
\end{lemma}

\begin{proof}
\leanok
A permutation is a bijection, so $|\sigma(P)| = |P|$. If $|P|\ne|Q|$
then $\sigma(P) \ne Q$ for all $\sigma$.
\end{proof}

\subsubsection{Expansion steps}

\begin{lemma}
\label{lem.imageFinset.compl}
\lean{AlgebraicCombinatorics.CauchyBinet.imageFinset_compl}
\leanhelper
\leanok
For a permutation $\sigma\in S_n$ and a subset $P\subseteq [n]$,
$\sigma(\overline{P}) = \overline{\sigma(P)}$.
\end{lemma}

\begin{proof}
\leanok
Since $\sigma$ is a bijection on $[n]$, the complement is preserved.
\end{proof}

\subsubsection{Permutation decomposition infrastructure}

The proof of Theorem~\ref{thm.det.det(A+B)} requires decomposing
permutations $\sigma \in S_n$ with $\sigma(P) = Q$ into independent
parts acting on~$P$ and on~$\overline{P}$.

\begin{definition}[Extraction of component permutations]
\label{def.extract}
\lean{AlgebraicCombinatorics.CauchyBinet.extractAlpha, AlgebraicCombinatorics.CauchyBinet.extractBeta}
\leanhelper
\leanok
Let $P,Q\subseteq [n]$ with $|P|=|Q|=k$ and let $\sigma\in S_n$
satisfy $\sigma(P) = Q$.  Write the elements of $P$ in increasing
order as $p_1 < \cdots < p_k$, of $Q$ as $q_1 < \cdots < q_k$,
of $\overline{P}$ as $p'_1 < \cdots < p'_\ell$, and of
$\overline{Q}$ as $q'_1 < \cdots < q'_\ell$.

Define $\alpha_\sigma \in S_k$ by
$\sigma(p_i) = q_{\alpha_\sigma(i)}$ and
$\beta_\sigma \in S_\ell$ by
$\sigma(p'_i) = q'_{\beta_\sigma(i)}$.
\end{definition}

\begin{definition}[Construction of permutation from components]
\label{def.constructSigma}
\lean{AlgebraicCombinatorics.CauchyBinet.constructSigma, AlgebraicCombinatorics.CauchyBinet.constructSigma_imageFinset, AlgebraicCombinatorics.CauchyBinet.constructSigma_mem_permsMapping}
\leanhelper
\leanok
Given $P,Q\subseteq [n]$ with $|P|=|Q|$ and permutations
$\alpha\in S_k$, $\beta\in S_\ell$ (with $k=|P|$, $\ell=|\overline{P}|$),
define $\sigma \in S_n$ by:
\[
\sigma(p_i) = q_{\alpha(i)} \quad\text{and}\quad
\sigma(p'_i) = q'_{\beta(i)}.
\]
This is the inverse of the map $\sigma \mapsto (\alpha_\sigma,\beta_\sigma)$.
\end{definition}

\begin{lemma}[Round-trip properties]
\label{lem.roundtrip}
\lean{AlgebraicCombinatorics.CauchyBinet.extractAlpha_constructSigma, AlgebraicCombinatorics.CauchyBinet.extractBeta_constructSigma, AlgebraicCombinatorics.CauchyBinet.constructSigma_extract}
\leanhelper
\leanok
The maps $\sigma \mapsto (\alpha_\sigma, \beta_\sigma)$ and
$(\alpha,\beta) \mapsto \sigma_{(\alpha,\beta)}$
are mutual inverses:
\begin{enumerate}
\item Extracting $\alpha$ from the permutation constructed from $(\alpha,\beta)$ returns $\alpha$.
\item Extracting $\beta$ from the permutation constructed from $(\alpha,\beta)$ returns $\beta$.
\item Constructing a permutation from the components $(\alpha_\sigma, \beta_\sigma)$ extracted from $\sigma$ returns $\sigma$.
\end{enumerate}
\end{lemma}

\begin{proof}
By unfolding the definitions and showing that the sorted-order
embeddings and their inverses compose to the identity.
\end{proof}

\begin{lemma}[Extract specification]
\label{lem.extract.spec}
\lean{AlgebraicCombinatorics.CauchyBinet.extractAlpha_spec, AlgebraicCombinatorics.CauchyBinet.extractBeta_spec}
\leanhelper
\leanok
For $\sigma$ with $\sigma(P) = Q$:
\begin{enumerate}
\item $\sigma(p_j) = q_{\alpha_\sigma(j)}$ for all $j$,
\item $\sigma(p'_j) = q'_{\beta_\sigma(j)}$ for all $j$.
\end{enumerate}
\end{lemma}

\begin{proof}
By definition of $\alpha_\sigma$ and $\beta_\sigma$, using the
order-embedding properties of sorted finsets.
\end{proof}

\begin{lemma}[Product factorization via extract]
\label{lem.prod.factor}
\lean{AlgebraicCombinatorics.CauchyBinet.prod_P_eq_submatrix_det, AlgebraicCombinatorics.CauchyBinet.prod_Pc_eq_submatrix_det}
\leanhelper
\leanok
For $\sigma$ with $\sigma(P) = Q$:
\begin{enumerate}
\item $\displaystyle\prod_{i\in P} A_{\sigma(i),i}
  = \prod_{j=1}^{k} A_{q_{\alpha(j)},\, p_j}$,
\item $\displaystyle\prod_{i\in\overline{P}} B_{\sigma(i),i}
  = \prod_{j=1}^{\ell} B_{q'_{\beta(j)},\, p'_j}$.
\end{enumerate}
\end{lemma}

\begin{proof}
Reindex the product over sorted elements of $P$ (resp.\ $\overline{P}$)
using Lemma~\ref{lem.extract.spec}.
\end{proof}

\subsubsection{Sign decomposition infrastructure}

The sign decomposition (Lemma~\ref{lem.sign.decomposition}) is proved
by reducing any $(P,Q)$ to the base case $P = Q = [k]$ via
``left shifts'' and ``left co-shifts.''

\begin{lemma}[Shift existence]
\label{lem.shift.exists}
\lean{AlgebraicCombinatorics.CauchyBinet.exists_shift_opportunity}
\leanhelper
\leanok
Let $0 < k < n$ and let $P\subseteq [n]$ with $|P| = k$ and
$P \ne [k]$.  Then there exists $i$ such that $i \notin P$ and
$i+1 \in P$ (with $i+1 < n$).
\end{lemma}

\begin{proof}
\leanok
If $P = [k]$ (the first $k$ elements), then $P$ is already the
prefix.  Otherwise, some element of $P$ is ``too large,'' meaning
there exists $j$ with $p_j > j$ (the $j$-th smallest element is
above position $j$).  A downward-closure argument produces the
required adjacent pair.
\end{proof}

\begin{lemma}[Sum decrement under shift]
\label{lem.shift.sum}
\lean{AlgebraicCombinatorics.CauchyBinet.finsetSumFin_swap_of_left_shift}
\leanhelper
\leanok
If $i\notin P$ and $i+1\in P$, let $P' = s_i(P)$ where
$s_i = \operatorname{swap}(i,i+1)$.  Then
$\operatorname{sum}(P') = \operatorname{sum}(P) - 1$.
\end{lemma}

\begin{proof}
\leanok
$P'$ is obtained from $P$ by replacing $i+1$ with $i$, so the sum
decreases by $(i+1)-i = 1$.
\end{proof}

\begin{lemma}[Combined sign preserved under left shift]
\label{lem.shift.combined.sign}
\lean{AlgebraicCombinatorics.CauchyBinet.leftShift_preserves_combined_sign}
\leanhelper
\leanok
If $i\notin P$ and $i+1\in P$, let $P' = s_i(P)$ and $\sigma' = \sigma\cdot s_i$.
Then
\[
(-1)^{\operatorname{sum} P' + \operatorname{sum} Q}\cdot (-1)^{\sigma'}
= (-1)^{\operatorname{sum} P + \operatorname{sum} Q}\cdot (-1)^{\sigma}.
\]
\end{lemma}

\begin{proof}
$(-1)^{\sigma'} = (-1)^{\sigma}\cdot(-1)^{s_i} = -(-1)^{\sigma}$
(since adjacent transpositions are odd), and
$(-1)^{\operatorname{sum} P'} = -(-1)^{\operatorname{sum} P}$
(since $\operatorname{sum} P' = \operatorname{sum} P - 1$ by
Lemma~\ref{lem.shift.sum}).  The two sign flips cancel.
\end{proof}

\begin{lemma}[Extract preserved under left shift]
\label{lem.shift.extract}
\lean{AlgebraicCombinatorics.CauchyBinet.extractAlpha_leftShift_eq, AlgebraicCombinatorics.CauchyBinet.extractBeta_leftShift_eq}
\leanhelper
\leanok
Under the left shift $P' = s_i(P)$, $\sigma' = \sigma\cdot s_i$:
\begin{enumerate}
\item $(-1)^{\alpha_{\sigma'}} = (-1)^{\alpha_\sigma}$,
\item $(-1)^{\beta_{\sigma'}} = (-1)^{\beta_\sigma}$.
\end{enumerate}
\end{lemma}

\begin{proof}
The sorted order of $P'$ is $s_i$ applied to the sorted order of~$P$
(since replacing $i{+}1$ by $i$ preserves relative order among all
other elements).  Therefore
$\sigma'(P'_j) = \sigma(s_i(s_i(P_j))) = \sigma(P_j)$,
so $\alpha_{\sigma'} = \alpha_\sigma$ (and similarly for $\beta$).
\end{proof}

\begin{lemma}[Left shift invariance]
\label{lem.shift.invariant}
\lean{AlgebraicCombinatorics.CauchyBinet.sign_decomposition_leftShift_invariant}
\leanhelper
\leanok
If the sign decomposition formula holds for $(P', Q, \sigma')$
(where $P' = s_i(P)$, $\sigma' = \sigma\cdot s_i$), then it
holds for $(P, Q, \sigma)$.
\end{lemma}

\begin{proof}
\leanok
Combine Lemma~\ref{lem.shift.combined.sign} (combined sign preserved)
with Lemma~\ref{lem.shift.extract} ($\alpha$ and $\beta$ signs preserved).
\end{proof}

\begin{lemma}[Left co-shift invariance]
\label{lem.coshift.invariant}
\lean{AlgebraicCombinatorics.CauchyBinet.sign_decomposition_leftCoShift_invariant}
\leanhelper
\leanok
If the sign decomposition formula holds for $(P, Q', \sigma')$
(where $Q' = s_i(Q)$, $\sigma' = s_i\cdot\sigma$ for $i\notin Q$,
$i{+}1\in Q$), then it holds for $(P, Q, \sigma)$.
\end{lemma}

\begin{proof}
\leanok
Analogous to Lemma~\ref{lem.shift.invariant}: the combined sign
is preserved under left co-shifts, and the extracted component
permutations $\alpha$ and $\beta$ are unchanged.
\end{proof}

\begin{lemma}[Sign decomposition base case]
\label{lem.sign.decomposition.prefix}
\lean{AlgebraicCombinatorics.CauchyBinet.sign_decomposition_prefix}
\leanhelper
\leanok
When $P = Q = [k]$ (the prefix finset), the sign decomposition formula
\[
(-1)^{\sigma}
= (-1)^{\operatorname{sum}[k]+\operatorname{sum}[k]}
\cdot (-1)^{\alpha_\sigma}\cdot (-1)^{\beta_\sigma}
= (-1)^{\alpha_\sigma}\cdot (-1)^{\beta_\sigma}
\]
holds, where the sign factor vanishes because
$\operatorname{sum}[k] + \operatorname{sum}[k] = k(k-1)$ is even.
\end{lemma}

\begin{proof}
\leanok
Since $\sigma$ preserves $[k]$, it decomposes as
$\sigma = \alpha \oplus \beta$
(a direct sum of the restrictions to $[k]$ and its complement).
The sign of a direct sum satisfies
$(-1)^\sigma = (-1)^\alpha \cdot (-1)^\beta$
(this is \texttt{Equiv.Perm.sign\_subtypeCongr} from Mathlib).
The factor $(-1)^{k(k-1)} = 1$ since $k(k-1)$ is even.
\end{proof}

\begin{lemma}
\label{lem.det.add.step1}
\lean{AlgebraicCombinatorics.CauchyBinet.det_add_expand_step1}
\leanhelper
\leanok
For $n\times n$-matrices $A$ and $B$,
\[
\det(A+B) = \sum_{\sigma\in S_n} (-1)^{\sigma}
\sum_{P\subseteq [n]}
\Bigl(\prod_{i\in P} A_{\sigma(i),i}\Bigr)
\Bigl(\prod_{i\in \overline{P}} B_{\sigma(i),i}\Bigr).
\]
\end{lemma}

\begin{proof}
Expand $\det(A+B) = \sum_{\sigma} (-1)^{\sigma} \prod_i (A_{\sigma(i),i} + B_{\sigma(i),i})$
and distribute each factor using the product rule for sums
(this is \texttt{Fintype.prod\_add} from Mathlib).
\end{proof}

\begin{lemma}
\label{lem.det.add.step2}
\lean{AlgebraicCombinatorics.CauchyBinet.det_add_expand_step2}
\leanhelper
\leanok
Swapping summation order:
\[
\det(A+B) = \sum_{P\subseteq [n]} \sum_{\sigma\in S_n}
(-1)^{\sigma}
\Bigl(\prod_{i\in P} A_{\sigma(i),i}\Bigr)
\Bigl(\prod_{i\in \overline{P}} B_{\sigma(i),i}\Bigr).
\]
\end{lemma}

\begin{proof}
By swapping the order of summation
(this is \texttt{Finset.sum\_comm} from Mathlib).
\end{proof}

\begin{lemma}
\label{lem.det.add.step3}
\lean{AlgebraicCombinatorics.CauchyBinet.det_add_expand_step3}
\leanhelper
\leanok
Partitioning by $Q = \sigma(P)$:
\[
\det(A+B) = \sum_{P\subseteq [n]} \sum_{Q\subseteq [n]}
\sum_{\substack{\sigma\in S_n;\\ \sigma(P) = Q}}
(-1)^{\sigma}
\Bigl(\prod_{i\in P} A_{\sigma(i),i}\Bigr)
\Bigl(\prod_{i\in \overline{P}} B_{\sigma(i),i}\Bigr).
\]
\end{lemma}

\begin{proof}
Split $\sum_{\sigma\in S_n}$ according to the value of
$Q = \sigma(P)$, using pairwise disjointness of the fibers.
\end{proof}

\begin{lemma}[Sign decomposition]
\label{lem.sign.decomposition}
\lean{AlgebraicCombinatorics.CauchyBinet.sign_decomposition}
\leanhelper
\leanok
Let $P,Q\subseteq [n]$ with $|P|=|Q|$ and let $\sigma\in S_n$ satisfy
$\sigma(P)=Q$. Write the elements of $P$ in increasing order as
$p_1<p_2<\cdots<p_k$, the elements of $Q$ as $q_1<q_2<\cdots<q_k$,
the elements of $\overline{P}$ as $p'_1<\cdots<p'_{\ell}$, and those
of $\overline{Q}$ as $q'_1<\cdots<q'_{\ell}$.

Define $\alpha_{\sigma}\in S_k$ by $\sigma(p_i) = q_{\alpha_{\sigma}(i)}$
and $\beta_{\sigma}\in S_{\ell}$ by
$\sigma(p'_i) = q'_{\beta_{\sigma}(i)}$.

Then
\[
(-1)^{\sigma} = (-1)^{\operatorname{sum} P + \operatorname{sum} Q}
\cdot (-1)^{\alpha_{\sigma}} \cdot (-1)^{\beta_{\sigma}}.
\]
\end{lemma}

\begin{proof}
The proof uses strong induction on
$\operatorname{sum} P + \operatorname{sum} Q$.

\textbf{Base case.}
When $P = Q = [k]$, the formula holds by
Lemma~\ref{lem.sign.decomposition.prefix}.

\textbf{Inductive case.}
If $P \ne [k]$, Lemma~\ref{lem.shift.exists} provides an index $i$
with $i\notin P$ and $i+1\in P$.  The left shift
$P' = s_i(P)$, $\sigma' = \sigma\cdot s_i$
reduces $\operatorname{sum} P$ by $1$, so the induction hypothesis
applies to $(P', Q, \sigma')$.
Lemma~\ref{lem.shift.invariant} lifts the result back to $(P, Q, \sigma)$.

Similarly, if $P = [k]$ but $Q \ne [k]$,
Lemma~\ref{lem.shift.exists} applied to $Q$ gives an index for a
left co-shift, and Lemma~\ref{lem.coshift.invariant} lifts the result.
\end{proof}

\begin{lemma}[Key factorization]
\label{lem.det.sumPQ}
\lean{AlgebraicCombinatorics.CauchyBinet.sum_perms_mapping_eq_det_product}
\leanhelper
\leanok
For fixed subsets $P,Q\subseteq [n]$ with $|P|=|Q|$:
\[
\sum_{\substack{\sigma\in S_n;\\ \sigma(P)=Q}}
(-1)^{\sigma}
\Bigl(\prod_{i\in P} A_{\sigma(i),i}\Bigr)
\Bigl(\prod_{i\in \overline{P}} B_{\sigma(i),i}\Bigr)
= (-1)^{\operatorname{sum} P + \operatorname{sum} Q}
\det(\operatorname{sub}_P^Q A)
\cdot \det(\operatorname{sub}_{\overline{P}}^{\overline{Q}} B).
\]
\end{lemma}

\begin{proof}
\leanok
The proof uses the bijection
$\sigma \mapsto (\alpha_{\sigma}, \beta_{\sigma})$
(Definition~\ref{def.extract})
from $\{\sigma\in S_n : \sigma(P) = Q\}$ to $S_k \times S_{\ell}$,
with inverse given by the construction in
Definition~\ref{def.constructSigma}.
The round-trip properties (Lemma~\ref{lem.roundtrip})
establish that this is a bijection.

After reindexing by $(\alpha,\beta)$:
\begin{enumerate}
\item The sign identity (Lemma~\ref{lem.sign.decomposition}) gives
  $(-1)^{\sigma} = (-1)^{\operatorname{sum} P+\operatorname{sum} Q}
  \cdot (-1)^{\alpha}\cdot (-1)^{\beta}$.
\item The products over $P$ and $\overline{P}$ factor
  (Lemma~\ref{lem.prod.factor}).
\item The double sum
  $\sum_{\alpha}\sum_{\beta}$ factors into a product of two Leibniz
  expansions, giving
  $\det(\operatorname{sub}_P^Q A)\cdot
   \det(\operatorname{sub}_{\overline{P}}^{\overline{Q}} B)$.
\end{enumerate}
\end{proof}

\begin{theorem}
\label{thm.det.det(A+B)}
\lean{AlgebraicCombinatorics.CauchyBinet.det_add_sum}
\leantarget
\leanok
Let $n\in\mathbb{N}$. For any subset $I$ of $[n]$, let
$\widetilde{I} = [n]\setminus I$ denote its complement.

Let $A$ and $B$ be two $n\times n$-matrices in $K^{n\times n}$. Then,
\[
\det(A+B) = \sum_{P\subseteq [n]}\;\;
\sum_{\substack{Q\subseteq [n];\\ |P|=|Q|}}
(-1)^{\operatorname{sum} P + \operatorname{sum} Q}
\det(\operatorname{sub}_P^Q A)
\cdot \det(\operatorname{sub}_{\widetilde{P}}^{\widetilde{Q}} B).
\]
\end{theorem}

\begin{proof}
The expansion from Lemma~\ref{lem.det.add.step3} gives a triple sum
$\sum_P \sum_Q \sum_{\sigma:\sigma(P)=Q}$.
Applying Lemma~\ref{lem.det.sumPQ} to each inner sum collapses the
$\sigma$-sum. When $|P| \ne |Q|$, the set of permutations mapping
$P$ to $Q$ is empty (Lemma~\ref{lem.permsMapping.empty}),
so the contribution is $0$.
\end{proof}

\subsection{Minors of diagonal matrices}

\begin{lemma}
\label{lem.det.minors-diag}
\lean{AlgebraicCombinatorics.CauchyBinet.det_diagonal_submatrix_eq, AlgebraicCombinatorics.CauchyBinet.det_diagonal_submatrix_off_diag}
\leantarget
\leanok
Let $n\in\mathbb{N}$. Let $d_1,d_2,\ldots,d_n\in K$.
Let
\[
D:=\operatorname{diag}(d_1,d_2,\ldots,d_n)
=\begin{pmatrix}
d_1 & 0 & \cdots & 0\\
0 & d_2 & \cdots & 0\\
\vdots & \vdots & \ddots & \vdots\\
0 & 0 & \cdots & d_n
\end{pmatrix}\in K^{n\times n}.
\]

\textbf{(a)} We have $\det(\operatorname{sub}_P^P D) = \prod_{i\in P} d_i$
for any subset $P$ of $[n]$.

\textbf{(b)} Let $P$ and $Q$ be two distinct subsets of $[n]$
satisfying $|P|=|Q|$. Then,
$\det(\operatorname{sub}_P^Q D) = 0$.
\end{lemma}

\begin{proof}
\textbf{(a)} The submatrix $\operatorname{sub}_P^P D$ is itself diagonal with entries
$d_i$ for $i\in P$. Thus its determinant is $\prod_{i\in P} d_i$.

The submatrix of a diagonal matrix via an order-preserving embedding
is again diagonal; then the determinant of a diagonal matrix is
the product of its diagonal entries
(this is \texttt{Matrix.det\_diagonal} from Mathlib).

\textbf{(b)} Since $P\ne Q$ but $|P|=|Q|$, there exists $i\in P\setminus Q$.
The row of $\operatorname{sub}_P^Q D$ corresponding to $i$ is all zeros
(since $D_{i,j}=0$ for $j\ne i$ and $i\notin Q$), so the determinant
is $0$ (by \texttt{Matrix.det\_eq\_zero\_of\_row\_eq\_zero} from Mathlib).
\end{proof}

\subsection{Determinant of $A+D$}

\begin{theorem}
\label{thm.det.det(A+D)}
\lean{AlgebraicCombinatorics.CauchyBinet.det_add_diagonal}
\leantarget
\leanok
Let $n\in\mathbb{N}$. Let $A$ and $D$ be two $n\times n$-matrices in
$K^{n\times n}$ such that the matrix $D$ is diagonal. Let
$d_1,d_2,\ldots,d_n$ be the diagonal entries of the diagonal matrix $D$. Then,
\[
\det(A+D) = \sum_{P\subseteq [n]}
\det(\operatorname{sub}_P^P A) \cdot \prod_{i\in [n]\setminus P} d_i.
\]
The minors $\det(\operatorname{sub}_P^P A)$ are called the
\emph{principal minors} of $A$.
\end{theorem}

\begin{proof}
Theorem~\ref{thm.det.det(A+B)} (applied to $B=D$) yields a double
sum. By Lemma~\ref{lem.det.minors-diag}\textbf{(b)}, the terms with
$P\ne Q$ vanish (since
$\det(\operatorname{sub}_{\widetilde{P}}^{\widetilde{Q}} D)=0$).
In the surviving terms ($P=Q$), the sign factor
$(-1)^{\operatorname{sum} P + \operatorname{sum} P} = 1$, and
Lemma~\ref{lem.det.minors-diag}\textbf{(a)} gives
$\det(\operatorname{sub}_{\widetilde{P}}^{\widetilde{P}} D)
= \prod_{i\in\widetilde{P}} d_i$.
\end{proof}

\subsection{Determinant of $x + D$}

\begin{lemma}
\label{lem.constPlusDiag.decomp}
\lean{AlgebraicCombinatorics.CauchyBinet.constPlusDiagMatrix_eq_add}
\leanhelper
\leanok
The matrix $F$ with entries $F_{i,j} = x + d_i\,[i{=}j]$ decomposes as
$F = A + D$ where $A$ is the $n\times n$ constant matrix with all
entries equal to~$x$ and $D = \operatorname{diag}(d_1,\ldots,d_n)$.
\end{lemma}

\begin{proof}
\leanok
Entry-wise verification.
\end{proof}

\begin{lemma}
\label{lem.det.const.matrix.singleton}
\lean{AlgebraicCombinatorics.CauchyBinet.det_submatrix_constMatrix_singleton}
\leanhelper
\leanok
For a singleton $P = \{p\}$, the $1\times 1$ submatrix of the
constant-$x$ matrix has determinant $x$.
\end{lemma}

\begin{proof}
\leanok
The $1\times 1$ matrix is $(x)$, which has determinant $x$
(this is \texttt{Matrix.det\_unique} from Mathlib).
\end{proof}

\begin{lemma}
\label{lem.det.const.matrix.zero}
\lean{AlgebraicCombinatorics.CauchyBinet.det_submatrix_constMatrix_eq_zero}
\leanhelper
\leanok
Let $A$ be the $n\times n$ constant matrix with all entries equal to $x$.
If $P\subseteq [n]$ has $|P| \ge 2$, then
$\det(\operatorname{sub}_P^P A) = 0$.
\end{lemma}

\begin{proof}
\leanok
When $|P|\ge 2$, the submatrix $\operatorname{sub}_P^P A$ is a constant
matrix of size $\ge 2$. All its rows are identical, so
the determinant is $0$.
\end{proof}

\begin{proposition}
\label{prop.det.x+ai}
\lean{AlgebraicCombinatorics.CauchyBinet.det_const_add_diagonal}
\leantarget
\leanok
Let $n\in\mathbb{N}$. Let $d_1,d_2,\ldots,d_n\in K$ and $x\in K$.
Let $F$ be the $n\times n$-matrix
\[
\begin{pmatrix}
x+d_1 & x & \cdots & x\\
x & x+d_2 & \cdots & x\\
\vdots & \vdots & \ddots & \vdots\\
x & x & \cdots & x+d_n
\end{pmatrix}\in K^{n\times n}.
\]
Then,
\[
\det F = d_1 d_2 \cdots d_n
+ x \sum_{i=1}^n d_1 d_2 \cdots \widehat{d_i} \cdots d_n,
\]
where the hat over ``$d_i$'' means ``omit the $d_i$ factor.''
\end{proposition}

\begin{proof}
Write $F = A + D$ (Lemma~\ref{lem.constPlusDiag.decomp}) where $A$ is the $n\times n$ constant matrix with
all entries $x$ and $D = \operatorname{diag}(d_1,\ldots,d_n)$.
By Theorem~\ref{thm.det.det(A+D)},
\[
\det(A+D) = \sum_{P\subseteq [n]}
\det(\operatorname{sub}_P^P A) \cdot \prod_{i\in [n]\setminus P} d_i.
\]
For $|P|\ge 2$, $\det(\operatorname{sub}_P^P A) = 0$
(Lemma~\ref{lem.det.const.matrix.zero}), so only $P=\varnothing$ and
singleton $P=\{p\}$ contribute:
\begin{align*}
\det F
&= \underbrace{\det(\operatorname{sub}_{\varnothing}^{\varnothing} A)}_{=1}
\cdot \underbrace{\prod_{i\in [n]} d_i}_{=d_1 d_2\cdots d_n}
+ \sum_{p=1}^{n}
\underbrace{\det(\operatorname{sub}_{\{p\}}^{\{p\}} A)}_{=x}
\cdot \prod_{i\in [n]\setminus\{p\}} d_i \\
&= d_1 d_2 \cdots d_n + x \sum_{i=1}^n
d_1 d_2 \cdots \widehat{d_i} \cdots d_n.
\end{align*}
\end{proof}

\subsection{Characteristic polynomial coefficients}

\begin{proposition}
\label{prop.det.charpol-explicit}
\lean{AlgebraicCombinatorics.CauchyBinet.det_charPoly_coeff, AlgebraicCombinatorics.CauchyBinet.det_charPoly_coeff'}
\leantarget
\leanok
Let $n\in\mathbb{N}$. Let $A\in K^{n\times n}$ be an $n\times n$-matrix.
Let $x\in K$. Let $I_n$ denote the $n\times n$ identity matrix. Then,
\[
\det(A + x I_n)
= \sum_{P\subseteq [n]}
\det(\operatorname{sub}_P^P A) \cdot x^{n-|P|}
= \sum_{k=0}^{n}
\biggl(\sum_{\substack{P\subseteq [n];\\ |P|=n-k}}
\det(\operatorname{sub}_P^P A)\biggr) x^k.
\]
\end{proposition}

\begin{proof}
The matrix $xI_n$ is diagonal with all entries equal to $x$.
Theorem~\ref{thm.det.det(A+D)} (applied to $D = xI_n$ and $d_i = x$)
yields
\[
\det(A+xI_n) = \sum_{P\subseteq [n]}
\det(\operatorname{sub}_P^P A) \cdot
\underbrace{\prod_{i\in [n]\setminus P} x}_{= x^{n-|P|}}.
\]
This is the first equality. The second follows by grouping the sum
according to the value of $|P|$ and substituting $n-k$ for $k$.

In the Lean formalization, $x I_n$ is written as
$x \cdot I_n$ and the diagonal form is obtained from the identity
$x \cdot I_n = \operatorname{diag}(x,x,\ldots,x)$.
The second form regroups by partitioning $\mathcal{P}([n])$ into
fibers by cardinality.
\end{proof}

\subsection{The Pascal matrix}

\begin{definition}[Pascal matrix and its LU factors]
\label{def.pascal}
\lean{AlgebraicCombinatorics.CauchyBinet.pascalMatrix, AlgebraicCombinatorics.CauchyBinet.pascalLowerTriangular, AlgebraicCombinatorics.CauchyBinet.pascalUpperTriangular}
\leanhelper
\leanok
Define the $n\times n$ Pascal matrix, the lower triangular factor $L$,
and the upper triangular factor $U$ by:
\begin{align*}
A_{i,j} &= \binom{i+j}{i}, &
L_{i,k} &= \binom{i}{k}, &
U_{k,j} &= \binom{j}{k},
\end{align*}
for $0\le i,j,k \le n-1$.
\end{definition}

\begin{lemma}[Triangularity and unit diagonal]
\label{lem.pascal.triangular}
\lean{AlgebraicCombinatorics.CauchyBinet.pascalLowerTriangular_is_lower, AlgebraicCombinatorics.CauchyBinet.pascalUpperTriangular_is_upper, AlgebraicCombinatorics.CauchyBinet.pascalLowerTriangular_diag, AlgebraicCombinatorics.CauchyBinet.pascalUpperTriangular_diag}
\leanhelper
\leanok
$L$ is lower triangular ($L_{i,k} = 0$ for $i < k$) with $L_{i,i} = 1$.
$U$ is upper triangular ($U_{k,j} = 0$ for $k > j$) with $U_{k,k} = 1$.
\end{lemma}

\begin{proof}
\leanok
$\binom{i}{k} = 0$ when $i < k$, and $\binom{i}{i} = 1$.
Similarly for $U$.
\end{proof}

\begin{lemma}
\label{lem.pascal.LU}
\lean{AlgebraicCombinatorics.CauchyBinet.pascal_eq_LU}
\leanhelper
\leanok
Then $A = L \cdot U$ (over $\mathbb{Z}$).
\end{lemma}

\begin{proof}
\leanok
For all $i,j$, we have (by Vandermonde's identity)
\[
A_{i,j} = \binom{i+j}{i}
= \sum_{k=0}^{n-1} \binom{i}{k} \binom{j}{k}
= \sum_{k=0}^{n-1} L_{i,k}\, U_{k,j}
= (LU)_{i,j}.
\]
The key combinatorial identity
$\sum_{k} \binom{m}{k}\binom{n}{k} = \binom{m+n}{n}$
(Vandermonde's identity) is used to establish the LU decomposition.
\end{proof}

\begin{lemma}
\label{lem.pascal.L.det}
\lean{AlgebraicCombinatorics.CauchyBinet.det_pascalLowerTriangular}
\leanhelper
\leanok
The lower triangular factor $L$ has $\det L = 1$.
\end{lemma}

\begin{proof}
\leanok
$L$ is lower triangular
($\binom{i}{k} = 0$ when $i < k$) and its diagonal entries are
$L_{i,i} = \binom{i}{i} = 1$
(Lemma~\ref{lem.pascal.triangular}).
The determinant of a triangular matrix is the product of its diagonal
entries, which is $1$.
\end{proof}

\begin{lemma}
\label{lem.pascal.U.det}
\lean{AlgebraicCombinatorics.CauchyBinet.det_pascalUpperTriangular}
\leanhelper
\leanok
The upper triangular factor $U$ has $\det U = 1$.
\end{lemma}

\begin{proof}
\leanok
$U$ is upper triangular
($\binom{j}{k} = 0$ when $k > j$) and its diagonal entries are
$U_{k,k} = \binom{k}{k} = 1$
(Lemma~\ref{lem.pascal.triangular}).
The determinant of a triangular matrix is the product of its diagonal
entries, which is $1$.
\end{proof}

\begin{proposition}
\label{prop.det.pascal-LU}
\lean{AlgebraicCombinatorics.CauchyBinet.det_pascal_matrix}
\leantarget
\leanok
Let $n\in\mathbb{N}$. Let $A$ be the $n\times n$-matrix
\[
\Bigl(\binom{i+j-2}{i-1}\Bigr)_{1\le i\le n,\ 1\le j\le n}
= \begin{pmatrix}
\binom{0}{0} & \binom{1}{0} & \cdots & \binom{n-1}{0}\\
\binom{1}{1} & \binom{2}{1} & \cdots & \binom{n}{1}\\
\vdots & \vdots & \ddots & \vdots\\
\binom{n-1}{n-1} & \binom{n}{n-1} & \cdots & \binom{2n-2}{n-1}
\end{pmatrix}.
\]
Then $\det A = 1$.
\end{proposition}

\begin{proof}
By Lemma~\ref{lem.pascal.LU}, $A = LU$. Therefore
\[
\det A = \det(LU) = \det L \cdot \det U = 1 \cdot 1 = 1,
\]
using the multiplicativity of the determinant
(Theorem~\ref{thm.det.CB.square}) and
Lemmas~\ref{lem.pascal.L.det} and~\ref{lem.pascal.U.det}.
\end{proof}
